% Grado 5 -- generado por tools/build_tex.py a partir de raw/grado_05.txt
\section{Grado 5}\label{grado:5}

% PDF p. 37
\dba{1}{Interpreta y utiliza los números naturales y racionales en su representación fraccionaria para formular y resolver problemas aditivos, multiplicativos y que involucren operaciones de potenciación.}

\evidencias

\begin{itemize}
  \item Interpreta la relación parte - todo y la representa por medio de fracciones, razones o cocientes.
  \item Interpreta y utiliza números naturales y racionales (fraccionarios) asociados con un contexto para solucionar problemas.
  \item Determina las operaciones suficientes y necesarias para solucionar diferentes tipos de problemas.
  \item Resuelve problemas que requieran reconocer un patrón de medida asociado a un número natural o a un racional (fraccionario).
\end{itemize}

\ejemplo

Don Marcos, el dueño de una finca productora de frutas y vegetales, ha decidido distribuir su lote para sembrar los productos que se muestran en la siguiente imagen.

\fig{g05_p37_1}{Cuadrícula dividida en regiones con frutas y objetos para expresar partes de un todo como fracciones.}

Expresa la fracción del total de la finca que representa cada una de las situaciones siguientes y justifica las respuestas y procedimientos empleados:

a) La porción de tierra que piensa utilizar Don Marcos para construir su casa. b) La porción de tierra que se utilizará para sembrar bananos. c) La porción de tierra que se utilizará para sembrar. d) La porción de tierra que no se utilizará para sembrar.


% PDF p. 37
\dba{2}{Describe y desarrolla estrategias (algoritmos, propiedades de las operaciones básicas y sus relaciones) para hacer estimaciones y cálculos al solucionar problemas de potenciación.}

\evidencias

\begin{itemize}
  \item Utiliza las propiedades de las operaciones con números naturales y racionales (fraccionarios) para justificar algunas estrategias de cálculo o estimación relacionados con áreas de cuadrados y volúmenes de cubos.
  \item Descompone un número en sus factores primos.
  \item Identifica y utiliza las propiedades de la potenciación para resolver problemas aritméticos.
  \item Determina y argumenta acerca de la validez o no de estrategias para calcular potencias.
\end{itemize}

\ejemplo

Un profesor representa el producto (3 2 x 2) x (2 2 x 3)en una hoja cuadriculada de la siguiente manera:

\fig{g05_p37_2}{Cálculo de 18 × 12 por descomposición: 2 × (9×9) + 6 × (3²) = 2 × (9²) + 6 × (3²) = 2 × 81 + 6 × 9 = 162 + 54 = 216.}

Al dividir de forma diferente la cuadrícula explora si es posible encontrar otra manera de representar el mismo producto.


% PDF p. 38
\dba{3}{Compara y ordena números fraccionarios a través de diversas interpretaciones, recursos y representaciones.}

\evidencias

\begin{itemize}
  \item Representa fracciones con la ayuda de la recta numérica.
  \item Determina criterios para ordenar fracciones y expresiones decimales de mayor a menor o viceversa.
\end{itemize}

\ejemplo

Camilo construyó tres cintas métricas de la misma longitud y dividió la unidad de cada una de ellas en diferentes partes. Luego representó una fracción en cada una de ellas, como se muestra a continuación.

\fig{g05_p38_1}{Regla graduada de 0 a 3 con una fracción señalada con flecha.}

\fig{g05_p38_2}{Regla graduada de 0 a 3 con la fracción 13/8 señalada.}

\fig{g05_p38_3}{Regla graduada de 0 a 3 con otra fracción señalada.}

Utiliza las cintas de Camilo e identifica si 9 es igual, mayor o menor a 11. 5			 6


% PDF p. 39
\dba{4}{Justifica relaciones entre superficie y volumen, respecto a dimensiones de figuras y sólidos, y elige las unidades apropiadas según el tipo de medición (directa e indirecta), los instrumentos y los procedimientos.}

\evidencias

\begin{itemize}
  \item Determina las medidas reales de una figura a partir de un registro gráfico (un plano).
  \item Mide superficies y longitudes utilizando diferentes estrategias (composición, recubrimiento, bordeado, cálculo).
  \item Construye y descompone figuras planas y sólidos a partir de medidas establecidas.
  \item Realiza estimaciones y mediciones con unidades apropiadas según sea longitud, área o volumen.
\end{itemize}

\ejemplo

Con una piola de 50 cm se hacen rectángulos diferentes. El perímetro de estos rectángulos es el mismo, determina si sus áreas permanecen iguales. Determina si se pueden hacer cajas de caras rectangulares de volúmenes diferentes pero en las que se necesite la misma cantidad de cartón para hacer sus moldes.


% PDF p. 39
\dba{5}{Explica las relaciones entre el perímetro y el área de diferentes figuras (variaciones en el perímetro no implican variaciones en el área y viceversa) a partir de mediciones, superposición de figuras, cálculo, entre otras.}

\evidencias

\begin{itemize}
  \item Compara diferentes figuras a partir de las medidas de sus lados.
  \item Calcula las medidas de los lados de una figura a partir de su área.
  \item Dibuja figuras planas cuando se dan las medidas de los lados.
  \item Propone estrategias para la solución de problemas relativos a la medida de la superficie de figuras planas.
  \item Reconoce que figuras con áreas diferentes pueden tener el mismo perímetro.
  \item Mide superficies y longitudes utilizando diferentes estrategias (composición, recubrimiento, bordeado, cálculo).
\end{itemize}

\ejemplo

Luisa y sus amigas quieren empacar unas tarjetas que tienen diferentes formas (triángulos y cuadriláteros) en sobres rectangulares. Antes de empacar las tarjetas, les ponen un hilo decorativo en todo el borde.

La cantidad de papel utilizado en las tarjetas es 126cm 2 , o 144cm 2 o 120cm 2 . Por ejemplo, una tarjeta en forma de triángulo rectángulo mide en sus lados perpendiculares 20 cm y 12 cm, otra en forma de cuadrado mide de lado 12 cm.

Determina otras dimensiones posibles para los lados de las tarjetas utilizando esas cantidades de papel. Además, la longitud de sus respectivos lados para establecer la cantidad de hilo que se emplea en cada tarjeta y discute acerca de la posibilidad de tener varias tarjetas de igual área pero diferente perímetro. Explica los procedimientos utilizados.


% PDF p. 40
\dba{6}{Identifica y describe propiedades que caracterizan un cuerpo en términos de la bidimensionalidad y la tridimensionalidad y resuelve problemas en relación con la composición y descomposición de las formas.}

\evidencias

\begin{itemize}
  \item Relaciona objetos tridimensionales y sus propiedades con sus respectivos desarrollos planos.
  \item Reconoce relaciones intra e interfigurales.
  \item Determina las mediciones reales de una figura a partir de un registro gráfico (un plano).
  \item Construye y descompone figuras planas y sólidos a partir de medidas establecidas.
  \item Utiliza transformaciones a figuras en el plano para describirlas y calcular sus medidas.
  \item Reconoce diferentes distribuciones de plantillas de un cuerpo en una superficie, las formas en que pueden acoplarse o encajar, lee la información que presenta la plantilla del cuerpo o su representación en un plano.
\end{itemize}

\ejemplo

La empresa Tortimax requiere un empaque para sus productos. El molde del empaque es el que se muestra en la figura y debe ser elaborado en cartón industrial. El tamaño de un pliego de este material es de 100 cm X 70 cm. Determina la cantidad de moldes del empaque que puede realizarse por pliego para aprovechar al máximo el material. Describe y discute acerca del procedimiento utilizado.


% PDF p. 41
\dba{7}{Resuelve y propone situaciones en las que es necesario describir y localizar la posición y la trayectoria de un objeto con referencia al plano cartesiano.}

\evidencias

\begin{itemize}
  \item Localiza puntos en un mapa a partir de coordenadas cartesianas.
  \item Interpreta los elementos de un sistema de referencia (ejes, cuadrantes, coordenadas).
  \item Grafica en el plano cartesiano la posición de un objeto usando direcciones cardinales (norte, sur, oriente y occidente).
  \item Emplea el plano cartesiano al plantear y resolver situaciones de localización.
  \item Representa en forma gráfica y simbólica la localización y trayectoria de un objeto.
\end{itemize}

\ejemplo

\fig{g05_p41_1}{Mapa urbano con marcadores de lugares para describir ubicaciones y rutas.}

Tatiana es una turista que ha venido a visitarnos. Ayuda a Tatiana a ubicarse a partir de un plano de la ciudad, municipio o barrio. Lo que prefieras:

\begin{itemize}
  \item Realiza un mapa a escala del barrio en un papel cuadriculado.
  \item Crea un sistema de referencia para que Tatiana pueda conocer los mejores lugares usando los puntos cardinales (Norte, Sur, Oriente, Occidente).
  \item Escribe un mensaje a Tatiana para indicar cómo realizar el recorrido.
  \item Da instrucciones para seguir una trayectoria que permita ir de un sitio a otro. Propone otras trayectorias posibles.
\end{itemize}

% PDF p. 41
\dba{8}{Describe e interpreta variaciones de dependencia entre cantidades y las representa por medio de gráficas.}

\evidencias

\begin{itemize}
  \item Propone patrones de comportamiento numéricos y patrones de comportamiento gráficos.
  \item Realiza cálculos numéricos, organiza la información en tablas, elabora representaciones gráficas y las interpreta.
  \item Trabaja sobre números desconocidos para dar respuestas a los problemas.
\end{itemize}

\ejemplo

\fig{g05_p41_2}{Vaso que se llena de agua bajo un grifo (capacidad y volumen).}

Un recipiente cilíndrico recto, se llena con una llave que vierte 4 litros de agua cada 2 minutos. El cilindro tiene capacidad de 28 litros. Determina cuánto tiempo tarda el recipiente cilíndrico en llenarse. Determina cuántos litros hay en el recipiente a los cinco minutos después de abrir la llave. Determina qué ocurre con el nivel del agua a los 16 minutos.


% PDF p. 42
\dba{9}{Utiliza operaciones no convencionales, encuentra propiedades y resuelve ecuaciones en donde están involucradas.}

\evidencias

\begin{itemize}
  \item Interpreta y opera con operaciones no convencionales.
  \item Explora y busca propiedades de tales operaciones.
  \item Compara las propiedades de las operaciones convencionales de suma, resta, producto y división con las propiedades de las operaciones no convencionales.
  \item Resuelve ecuaciones numéricas cuando se involucran operaciones no convencionales.
\end{itemize}

\ejemplo

Representa verbales mediante expresiones numéricas: la multiplicación entre la suma de 24 más 45, y la resta de 24 menos 12. El doble de un número; el doble de un número aumentado en 5. La mitad de un número, la tercera parte de un número. Resuelve la ecuación: el doble de un número más 3 es igual a 9, encuentra el número.


% PDF p. 42
\dba{10}{Formula preguntas que requieren comparar dos grupos de datos, para lo cual recolecta, organiza y usa tablas de frecuencia, gráficos de barras, circulares, de línea, entre otros. Analiza la información presentada y comunica los resultados.}

\evidencias

\begin{itemize}
  \item Formula preguntas y elabora encuestas para obtener los datos requeridos e identifica quiénes deben responder.
  \item Registra, organiza y presenta la información recolectada usando tablas, gráficos de barras, gráficos de línea, y gráficos circulares.
  \item Selecciona los gráficos teniendo en cuenta el tipo de datos que se va a representar.
  \item Interpreta la información obtenida y produce conclusiones que le permiten comparar dos grupos de datos de una misma población.
  \item Escribe informes sencillos en los que compara la distribución de dos grupos de datos.
\end{itemize}

\ejemplo

\fig{g05_p42_1}{Representaciones de datos: diagrama de barras y diagrama circular, con zapatos como tema.}

La alcaldía del municipio ha programado una serie de actividades deportivas y recreativas, y ha solicitado al colegio un informe en el que se indique las tallas de los zapatos y de los uniformes de los estudiantes de cuarto y quinto de primaria. Para dar solución a la situación, elabora una encuesta, recolecta la información y redacta un informe con los resultados obtenidos en el cual incluye tablas y gráficos y analiza la información por cursos, por género y el comportamiento general de los dos grados.


% PDF p. 43
\dba{11}{Utiliza la media y la mediana para resolver problemas en los que se requiere presentar o resumir el comportamiento de un conjunto de datos.}

\evidencias

\begin{itemize}
  \item Interpreta y encuentra la media y la mediana en un conjunto de datos usando estrategias gráficas y numéricas.
  \item Explica la información que brinda cada medida en relación con el conjunto de datos.
  \item Selecciona una de las medidas como la más representativa del comportamiento del conjunto de datos estudiado.
  \item Argumenta la selección realizada empleando semejanzas y diferencias entre lo que cada una de las medidas indica.
\end{itemize}

\ejemplo

Una campaña emprendida por el Ministerio de Salud y Protección Social para prevenir el aumento en los índices de obesidad y diabetes infantil y juvenil, sugiere que en promedio cada persona debe realizar 30 minutos diarios de una actividad física aeróbica de intensidad moderada (caminar, trotar, correr, nadar, montar en bicicleta, etc.), para evitar el sobrepeso. Se afirma que para cumplir con la campaña, cada persona debe hacer exactamente 30 minutos de ejercicio diarios. Argumenta la validez de esta afirmación.


% PDF p. 43
\dba{12}{Predice la posibilidad de ocurrencia de un evento simple a partir de la relación entre los elementos del espacio muestral y los elementos del evento definido.}

\evidencias

\begin{itemize}
  \item Reconoce situaciones aleatorias en contextos cotidianos.
  \item Enumera todos los posibles resultados de un experimento aleatorio simple.
  \item Identifica y enumera los resultados favorables de ocurrencia de un evento simple.
  \item Anticipa la ocurrencia de un evento simple.
\end{itemize}

\ejemplo

En un día de la recreación se realizan diferentes actividades con juegos de azar. Javier y Arturo eligen el juego de la ruleta. Las reglas acordadas son:

\begin{itemize}
  \item Cada uno selecciona una ruleta (Ruleta 1 o Ruleta 2).
  \item Al mismo tiempo giran una vez cada ruleta.
  \item Javier gana si saca un número par.
  \item Arturo gana si saca un número impar.
  \item Si Javier saca impar y Arturo saca par, vuelven a jugar.
\end{itemize}
\fig{g05_p44_1}{Ruleta de 12 sectores numerados (algunos números en rojo).}

\fig{g05_p44_2}{Ruleta de 20 sectores numerados.}

Reconoce que el juego de la ruleta corresponde a una situación aleatoria, identifica los eventos, asigna la probabilidad de ocurrencia y da argumentos para decidir si el juego es o no justo estadísticamente.


